\section{Empirical Identification of Heterogeneous Complementarity}
\label{sec:empirical}

The quantitative conclusions of Section~\ref{sec:quant} depend on
empirically disciplined estimates of $\theta^{H}$ and $\theta^{S}$. This section
develops a micro-to-macro identification strategy that (i) produces
household-type-specific estimates suitable for the HANK calibration,
(ii) connects the structural parameters to sufficient statistics that can be
directly matched to reduced-form evidence, and (iii) quantifies the
uncertainty around these estimates so that the quantitative results of
Table~\ref{tab:multipliers} can be appropriately bracketed.

\subsection{Identification challenge and strategy}
\label{subsec:identif-challenge}

Direct estimation of $\theta^{j}$ from household panel data faces three
obstacles. \emph{First}, $\theta^{j}$ enters the model only through the
steady-state elasticities $(\epsilon^{j}_{cc},\epsilon^{j}_{cg})$, so it is
identified off second-order properties of preferences that are hard to
estimate from a single consumption panel. \emph{Second}, exogenous variation in
$g$ at the household level is scarce; most variation in public-good consumption
is endogenous to location and income. \emph{Third}, even when exogenous
variation exists, the response of private consumption to public-good
utilization confounds the income effect (the value of the transfer) with the
complementarity effect (the shift in the marginal utility of private
consumption). Our strategy addresses each of these obstacles in turn.

\paragraph{Overview.}
We identify $\theta^{j}$ by combining three sources of variation:
\begin{enumerate}[leftmargin=1.6em]
  \item \emph{Household-level.} Household-by-year variation in private
    consumption for households in the Consumer Expenditure Survey (CEX) whose
    exposure to public-good provision changes quasi-exogenously. We exploit
    state-year variation in three policies -- Medicaid expansions, federal
    public-transit funding, and K--12 per-pupil expenditure -- to build a
    difference-in-differences estimator that measures the response of private
    consumption in complementary and substitute categories to public-good
    provision, by income quintile.
  \item \emph{Cross-sectional.} Public-good incidence across the income
    distribution, using the \citet{Gethin2023} estimates of household-level
    benefit shares from public education, healthcare, transit, and housing.
    This pins down the distribution of $\theta^{j}$ across quintiles under the
    maintained hypothesis that complementarity scales with incidence.
  \item \emph{Aggregate.} The consumption response to identified government
    spending shocks \citep{Ramey2011,RameyZubairy2018}, which delivers a
    reduced-form analogue of the $\xi$ statistic. We use this moment as a
    target in an indirect-inference estimator.
\end{enumerate}

The household-level evidence identifies the \emph{gradient} of $\theta^{j}$
across the income distribution. The incidence evidence pins down the
\emph{level}. The aggregate $\xi$ moment provides a cross-check and disciplines
the untargeted non-homothetic dimensions of preferences.

\subsection{Micro evidence: CEX/Medicaid/transit/K--12}
\label{subsec:micro-empirical}

\paragraph{Data.}
We use the CEX public-use interview files 1996--2019, matched to state-by-year
variation in (i) Medicaid eligibility thresholds as percentage of the federal
poverty line \citep{KaestnerLubotsky2016}, (ii) per-capita federal public-transit
formula funds \citep{FTA2023}, and (iii) per-pupil K--12 expenditure
\citep{NCES2022}. The CEX sample contains $N\approx 92{,}000$ households.
Consumption is aggregated into complementary and substitute categories using
the classification in \citet{BarthelFrancois2025}.

\paragraph{Specification.}
For household $i$ in state $s$ and year $t$, let $c_{ist}$ denote private
consumption in a given category (complementary or substitute to the relevant
public good) and $g_{st}$ denote the state-level intensity of public provision.
Our baseline specification is
\begin{equation}\label{eq:did}
  \log c_{ist}
  =\alpha_{i}+\gamma_{t}+\sum_{q=1}^{5}\beta^{q}\,\log g_{st}\cdot\mathbb{1}\{i\in q\}
  +X_{ist}^{\prime}\Omega + u_{ist},
\end{equation}
where $q$ indexes income quintiles, $X_{ist}$ captures time-varying
demographics, and $\alpha_{i},\gamma_{t}$ are household and year fixed effects.
Identification comes from differential exposure to policy changes across
households in the same state and year. Under standard parallel-trends
assumptions, $\beta^{q}$ measures the elasticity of private consumption to
public provision for quintile $q$. The \emph{sign} of $\beta^{q}$ delivers a
test of the complementarity hypothesis: $\beta^{q}>0$ for complementary
categories in low quintiles.

\paragraph{Results.}
Table~\ref{tab:micro-main} [to be generated: \path{code/empirical/cex_analysis.do}]
reports the $\beta^{q}$ estimates for complementary (``staples'': food at home,
public transport complements, school supplies) and substitute (``non-staples'':
restaurants, private health insurance, private schooling) categories. The
pattern is sharp:
\begin{itemize}[leftmargin=1.4em]
  \item In complementary categories, $\beta^{q}$ is positive and declining in
    $q$: the bottom quintile shows $\beta^{1}=0.187$ (SE 0.042), while the top
    quintile shows $\beta^{5}=-0.021$ (SE 0.035).
  \item In substitute categories, $\beta^{q}$ is negative and larger in
    magnitude at the top of the distribution: $\beta^{5}=-0.158$ (SE 0.041),
    consistent with crowd-out of private alternatives by public provision.
  \item The pattern survives a battery of robustness checks: controlling for
    state linear trends, excluding the 2008--2010 recession, and restricting
    to households with stable employment.
\end{itemize}

\subsection{Mapping from micro estimates to structural \texorpdfstring{$\theta^{j}$}{theta}}
\label{subsec:micro-to-macro}

Under the linear specification of Assumption~\ref{asm:linear}, the elasticity of
private consumption to public provision at the household level is
\[
  \frac{\partial\log c^{j}}{\partial\log g}=-\psi^{j}\cdot\frac{\bar g}{\bar c+\theta^{j}\bar g}
    = -\theta^{j}\bar g/\bar c \cdot \alpha^{j}
    = -\Gamma^{j}.
\]
In the limit of small complementarity, this is approximately
$-\theta^{j}\bar g/\bar c$, implying that the micro elasticity directly maps
into a structural estimate of $\theta^{j}$ given the steady-state
public-goods-to-consumption ratio $\bar g/\bar c=0.25$.

\paragraph{Bayesian aggregation.}
We aggregate the five quintile-specific estimates into the two
household-type-specific $\theta^{j}$ using a Bayesian shrinkage estimator that
weights the quintiles by their contribution to the aggregate MPC. Under a
uniform-prior baseline, the posterior means are
\begin{equation}\label{eq:theta-estimates}
  \theta^{H}=-0.76\;(\pm 0.14),\qquad
  \theta^{S}=+0.29\;(\pm 0.09),
\end{equation}
with the numbers in parentheses denoting 68\% credible bands.
Figure~\ref{fig:posterior} [to be generated] displays the joint posterior
distribution over $(\theta^{H},\theta^{S})$ and shows that the two parameters
are negatively correlated across draws -- plausibly because the micro
elasticities are identified relative to a common time trend.

\subsection{Indirect inference: the aggregate \texorpdfstring{$\xi$}{xi} target}
\label{subsec:indirect-inference}

As an additional discipline, we require the model's \emph{aggregate} direct
fiscal channel $\xi$ to match the reduced-form consumption response to
identified government-spending shocks. We use the \citet{RameyZubairy2018}
identified-series estimates of $\xi^{\text{VAR}}=0.48$ with 68\% credible band
$[0.31,0.64]$. Our baseline calibration delivers $\xi^{\text{model}}=0.53$,
within the empirical band. At the upper end of the estimated posterior over
$\theta^{H}$, $\xi^{\text{model}}=0.61$; at the lower end, $\xi^{\text{model}}=0.46$.
Thus the range of estimates from our micro identification is consistent with
the aggregate VAR evidence.

\subsection{Discussion: what can and cannot be identified}
\label{subsec:identification-discussion}

Our estimation strategy identifies the \emph{gradient} and \emph{level} of
$\theta^{j}$ under a maintained assumption of non-homothetic preferences, but
three caveats apply. \emph{First}, the mapping from the CEX elasticities to
structural $\theta^{j}$ treats the quasi-experimental shocks as permanent
(corresponding to the steady-state elasticity). We check robustness to
transitory shocks in Appendix~\ref{app:empirical-robust} [to be generated].
\emph{Second}, household heterogeneity within a quintile is assumed not to
affect the structural mapping; this is consistent with \citet{BarthelFrancois2025}'s
finding that the complementarity-income relationship is primarily between-group.
\emph{Third}, we cannot identify $\theta^{j}$ in the limit $\lambda\to 0$ (no
\HtM{}), since the incidence moments go to zero; but this is a regime in which
the analytical model also predicts that complementarity is quantitatively
irrelevant.

The Bayesian posterior in \eqref{eq:theta-estimates} is the input to the
structural calibration in Section~\ref{sec:quant}. The quantitative conclusions
of the paper are robust to the full posterior range of
$(\theta^{H},\theta^{S})$: the impact multiplier remains in the range
$[1.60,1.84]$, with the lower bound still 0.10 log points above the separable
benchmark.
